\section{Границы исследования и статус реализации}

Методология данной предварительной дипломной работы имеет два уровня. Первый уровень --- реализованный пайплайн данных и оценки, уже присутствующий в репозитории. Второй уровень --- предлагаемая обучающая цель, а именно multi-reference verifier-free (MRVF) reinforcement learning, которая математически сформулирована, но еще не полностью реализована в коде. Это различие важно для границ настоящей работы. Репозиторий уже содержит пайплайны для построения корпуса, генерации embeddings, извлечения ключевых слов, retrieval референсов, генерации кандидатов и автоматической pairwise evaluation. Однако сам этап обучения остается будущей работой. Поэтому данная глава различает то, что уже операционально в кодовой базе, и то, что предлагается как следующий этап проекта.

На высоком уровне реализованная система строит weakly supervised multi-reference датасет из сырых датасетов шуток. Сначала шутки объединяются в единый корпус, затем кодируются dense retrieval моделью и сводятся к небольшому набору ключевых слов. Эти ключевые слова расширяются в n-gram prompt groups, которые затем используются для retrieval семантически близких шуток из embedding index. Полученные референсы имеют два непосредственных применения. Во-первых, они дают данные для baseline generation и автоматической оценки. Во-вторых, они задают training и validation sets, необходимые для предлагаемой MRVF training objective.

\section{Построение датасета}

Датасет объединяется из двух публичных коллекций шуток: Short Jokes dataset и rJokes dataset. В репозитории этот этап реализован в \texttt{src/pipelines/jokes.py}. Пайплайн скачивает сырые файлы из зафиксированных revisions источников, преобразует их в parquet формат и сохраняет metadata источника для каждой шутки. Для каждой строки сохраняются поля \texttt{id}, \texttt{text}, \texttt{source\_name}, \texttt{source\_filename} и \texttt{source\_id}. После отдельной предобработки двух корпусов они конкатенируются и получают новый глобальный целочисленный идентификатор. Итоговый датасет содержит примерно \(664\,\text{k}\) объектов.

Такой дизайн намеренно является консервативным. Пайплайн не выполняет агрессивную нормализацию текста, кроме удаления пустых записей и стандартизации формата хранения. Для юмора такая осторожность желательна, поскольку кажущиеся мелкими текстовыми детали, такие как пунктуация, формулировка или переносы строк, могут влиять на комический эффект. Сохранение source metadata также полезно для последующего анализа доменных различий, поиска дубликатов и возможных source-specific biases.

\section{Dense embedding index}

Следующий этап кодирует каждую шутку в плотном семантическом векторном пространстве. Он реализован в \texttt{src/pipelines/embeddings.py}. Текущая default configuration использует модель \texttt{Qwen3-Embedding-8B} с размерностью embedding 1024, уменьшенной с исходных 4096 с помощью Matryoshka Representation Learning. Embeddings вычисляются асинхронно батчами, записываются в parquet shards и хранятся как пары \((i, e_i)\), где \(i\) --- глобальный идентификатор шутки, а \(e_i \in \mathbb{R}^{1024}\) --- соответствующий embedding.

Роль этого этапа двоякая. Во-первых, embeddings шуток затем используются для скоринга candidate keywords по семантической релевантности. Во-вторых, они задают vector index, используемый для retrieval multi-reference joke sets. Иными словами, embedding space является общим представлением, от которого зависят и keyword extraction, и reference retrieval.

\section{Извлечение ключевых слов как weak prompt retrieval}

Предполагается, что каждая шутка имеет набор weak prompts, на которые она правдоподобно отвечает. Например, ``write a joke about a chicken'' рассматривается как допустимый prompt для шутки ``why did the chicken crossed the road''. Этап keyword extraction направлен на построение weak prompts на основе ключевых слов шуток. Он реализован в \texttt{src/pipelines/keywords.py}.

Для каждой шутки \(y_i\) пайплайн сначала генерирует пул candidate n-grams с помощью analyzer \texttt{CountVectorizer}. В default configuration этот analyzer использует только unigrams, удаляет английские stopwords и оставляет не более 64 candidate terms на шутку. Пусть полученный candidate set равен
\[
    C(y_i)=\{c_{i1},\dots,c_{im}\}.
\]
Затем каждый candidate term кодируется той же embedding model, что и шутки. Если \(e(y_i)\) обозначает embedding шутки, а \(e(c)\) обозначает embedding candidate keyword, то relevance score вычисляется через cosine similarity:
\[
    s(c,y_i)=\frac{e(c)^\top e(y_i)}{\lVert e(c)\rVert\,\lVert e(y_i)\rVert}.
\]

Выбор только наиболее похожих candidates часто давал бы избыточные ключевые слова. Для уменьшения redundancy пайплайн применяет maximal marginal relevance (MMR), выбирая несколько ключевых слов с балансом между relevance и diversity. В default configuration diversity coefficient равен 0.7, а число сохраняемых keywords равно \(3\). Поэтому выход этого этапа --- компактное keyword representation
\[
    K(y_i)=\{k_{i1},k_{i2},k_{i3}\},
\]
возможно содержащее меньше трех элементов, если candidate pool мал.

Методологически эти ключевые слова не следует интерпретировать как полноценные латентные prompts. Это более слабые объекты: n-grams, которые захватывают тематические или лексические якоря шутки. Это прагматическое приближение. Оно не решает joke-to-prompt inversion полностью, но дает удобное представление, из которого можно построить несколько prompt options.

\section{Построение референсов через группировку ключевых слов и retrieval}

\subsection{Группировка ключевых слов и рендеринг prompts}

Этап reference construction реализован в \texttt{src/pipelines/references.py}. Для каждого keyword set \(K(y_i)\) пайплайн перечисляет все уникальные комбинации keywords, размер которых лежит между \texttt{min\_keywords} и \texttt{max\_keywords}. В default configuration эти значения равны 1 и 2, поэтому пайплайн строит все однословные и двухсловные группы. Если
\[
    K(y_i)=\{k_1,k_2,k_3\},
\]
то generated groups равны
\[
    \{k_1\},\ \{k_2\},\ \{k_3\},\ \{k_1,k_2\},\ \{k_1,k_3\},\ \{k_2,k_3\}.
\]
Каждая группа преобразуется в natural language prompt с помощью шаблона
\begin{quote}
    \ttfamily
    Write a joke using the following keyword(s): \{keywords\}
\end{quote}
Этот prompt служит двум целям: это текстовая форма, которая затем подается generation models, и это же объект, который кодируется как query для reference retrieval.

Такой дизайн создает контролируемый, но разнообразный спектр prompts. Single keyword prompts являются широкими, тематическими и часто допускают много возможных шуток. Two keyword prompts более ограничены и обычно извлекают более точные candidates. Текущая реализация останавливается на двух keywords, потому что более крупные группы быстро становятся слишком специфичными, уменьшая число правдоподобных различных референсов на prompt.

\subsection{Reference retrieval с embedding index}

После формирования prompt strings пайплайн кодирует их как queries и извлекает близкие шутки из FAISS index, построенного по полной коллекции joke embeddings. Retrieval index является inverted file index (IVF). Перед indexing и search векторы нормализуются, так что inner product соответствует cosine similarity. В default configuration основные параметры таковы:
\begin{itemize}
    \item embedding model: \texttt{Qwen3-Embedding-8B};
    \item dimensionality: 1024;
    \item \texttt{faiss\_nlist}: 4096;
    \item \texttt{faiss\_nprobe}: 32;
    \item \texttt{top\_k}: 5;
    \item \texttt{oversample}: 10;
    \item \texttt{min\_similarity}: 0.5.
\end{itemize}

Для rendered prompt \(x\) пусть \(q(x)\) --- его normalized embedding, а \(e_j\) --- normalized embedding шутки \(y_j\). Тогда retrieval score равен
\[
    s(x,y_j)=q(x)^\top e_j.
\]
Процедура поиска извлекает больше чем \(k\) candidates, фильтрует те, чья similarity ниже threshold, и оставляет top \(k\) remaining items. Соответствующие тексты затем используются как observed reference set:
\[
    Y^{*}(x)=\{y_1,\dots,y_n\}, \qquad n \le 5.
\]

На практике reference retrieval дает около \(4\) референсов на keyword group после deduplication и splitting. Target sample size для обучения равен \(10-20\text{k}\) references, поэтому достаточно сгенерировать \(3-5\text{k}\) keyword groups.

Этот этап реализует центральную идею работы: prompt связывается с несколькими candidate joke references, а не с одним ground truth. Важно, что retrieved set не является исчерпывающе размеченным множеством всех хороших шуток для \(x\). Это retrieved approximation такого множества. Это существенно и эмпирически, и теоретически. Эмпирически retrieved references могут содержать шум или пропускать валидные альтернативы. Теоретически неполнота \(Y^{*}(x)\) становится ключевой частью дальнейшего анализа MRVF, а не второстепенным неудобством.

\subsection{Deduplication и splitting}

После retrieval итоговый датасет содержит строки вида
\[
    (x_i, Y^{*}(x_i), S(x_i)),
\]
где \(x_i\) --- keyword-based prompt, \(Y^{*}(x_i)\) --- retrieved reference set, а \(S(x_i)\) --- vector of retrieval scores. Разные source jokes могут генерировать одинаковые keyword groups, поэтому репозиторий дедуплицирует датасет по полю \texttt{keywords}, сохраняет первое occurrence и заново назначает row identifiers.

Дедуплицированный датасет затем разбивается на train, validation и test partitions. Текущая реализация использует two-stage random split с \texttt{test\_fraction}=0.1, \texttt{validation\_fraction}=0.1 и fixed random seed. Отдельные parquet shards записываются для полного датасета и для каждого split.

\section{Генерация кандидатов}

Этот этап поддерживает candidate generation поверх reference dataset и реализован в \texttt{src/pipelines/candidates.py}. Для выбранной модели и split датасета пайплайн проходит по prompts и вызывает OpenAI-compatible chat completion endpoint. Используется тот же prompt template, что и при построении референсов:
\begin{quote}
    \ttfamily
    Write a joke using the following keyword(s): \{keywords\}
\end{quote}
В default configuration candidate generation использует temperature \(1.0\) и допускает до 128 completion tokens. Важно сохранять шутки короткими, чтобы упростить evaluation как в LLM-as-a-judge, так и в human settings.

Output schema содержит четыре поля: row identifier, keyword group, model name и generated text. Candidate files сохраняются в directory structure по model и split. Такая организация важна для последующей evaluation, поскольку позволяет выравнивать outputs нескольких моделей по одному underlying identifier и prompt.

Методологически этот candidate stage следует понимать как benchmarking layer. Он создает artifacts, необходимые для сравнения baseline models: pre-trained base, instruct и thinking versions, а также будущих trained checkpoints.

\section{Оценка кандидатов}

\subsection{Pairwise comparison}

Автоматическая evaluation реализована в \texttt{src/pipelines/evaluation.py}. Пайплайн сначала конкатенирует candidate datasets от нескольких моделей. Затем он группирует candidates по reference identifier и model name и строит pairwise comparisons для всех пар моделей, которые произвели outputs для одного prompt. Left-right ordering рандомизируется с fixed seed, чтобы уменьшить position bias.

Evaluation prompt намеренно прост: он содержит общий joke prompt, а также left и right candidate. System instruction просит annotator действовать как judge of humor quality, выбрать, какая шутка смешнее для general audience, избегать ties и возвращать strict JSON. В default configuration judge model работает при temperature \(0\).

Этот evaluation layer захватывает только comparative funniness при общем prompt. Он пока не реализует multidimensional evaluation, обсуждаемую в обзоре литературы: отдельные rubric-based judgments для prompt adherence, novelty, human-likeness и funniness; также он пока не включает human annotation. Эти этапы остаются future work, поскольку зависят от завершенной training procedure.

\subsection{Leaderboard aggregation}

Pairwise judgments агрегируются в leaderboard с wins, losses, числом comparisons, win rate и Bradley-Terry scores. Пусть \(\mathcal{M}\) обозначает множество evaluated models. Для каждой judged pair \((m_a,m_b)\) evaluator записывает binary win для одной стороны. Затем пайплайн оценивает Bradley-Terry strengths и сообщает их логарифмы как значения \texttt{bt\_score}.

Такая агрегация разумна для юмора, потому что relative judgments часто стабильнее absolute scores. Судье обычно легче решить, какая из двух шуток смешнее, чем дать хорошо калиброванную числовую оценку. В то же время эта процедура наследует ограничения judge model. Поскольку humor judgments языковых моделей могут расходиться с человеческими суждениями, текущий leaderboard следует интерпретировать как быстрый и дешевый proxy для intermediate validation, а не как окончательную меру качества юмора.

\section{Теоретическая рамка MRVF training}

Реализованный выше пайплайн строит prompt-conditioned reference sets; данный раздел определяет, как такие множества входят в reasoning-oriented reinforcement-learning objective. Центральное утверждение состоит в том, что генерация юмора плохо согласуется с single-reference assumptions, используемыми в стандартном verifier-based или verifier-free reasoning RL. Weak prompt может допускать множество приемлемых шуток, и репозиторий уже делает это конкретным, извлекая несколько референсов для одного keyword-based prompt. Ниже формализуется оптимизация по такой set-valued supervision.

\subsection{Обозначения}

Пусть \(\mathcal{X}\), \(\mathcal{Z}\) и \(\mathcal{Y}\) обозначают пространства prompts, reasoning traces и final answers соответственно. Для prompt \(x\in\mathcal{X}\):
\begin{itemize}
    \item \(z\in\mathcal{Z}\) обозначает reasoning trace,
    \item \(y\in\mathcal{Y}\) обозначает final generated joke,
    \item \(y^{*}\in\mathcal{Y}\) обозначает single reference answer,
    \item \(Y^{*}(x)\subset\mathcal{Y}\) обозначает \emph{observed} set of acceptable references, связанный с prompt \(x\),
    \item \(\pi_{\theta}\) обозначает policy, то есть условную вероятностную модель над текстом, параметризованную \(\theta\).
\end{itemize}

Более явно,
\[
    \pi_{\theta}(z,y\mid x)=P_{\theta}(Z=z,Y=y\mid X=x).
\]
Policy факторизуется как
\[
    \pi_{\theta}(z,y\mid x)=\pi_{\theta}(z\mid x)\,\pi_{\theta}(y\mid x,z).
\]
Когда модель явно выводит reasoning tokens, \(z\) можно отождествить с текстом внутри специального сегмента \texttt{<think>}, а \(y\) --- с final answer, который следует после него. Когда reasoning явно не экспонируется в output stream, ту же факторизацию можно понимать как концептуальное разложение на sampled intermediate trajectory \(z\) и final answer, conditioned on that trajectory. Преимущество такой нотации в том, что она разделяет две роли модели: выбор reasoning path и распределение probability mass по answers при данном path.

Для краткости optimization objective ниже записывается через observed set \(Y^{*}(x)\). При обсуждении incompleteness будет полезно представлять более крупное latent set \(Y_{\mathrm{true}}(x)\supseteq Y^{*}(x)\), содержащее все действительно приемлемые шутки для prompt \(x\), хотя это полное множество никогда не наблюдается на практике.

\subsection{RL с верифицируемыми reward}

Reinforcement learning оптимизирует policy, максимизируя expected reward. При заданной objective \(J(\theta)\) обучение идет через stochastic gradient ascent или, эквивалентно, через gradient descent по отрицательной objective. Поэтому критический design choice --- reward function: она определяет, какие generations получают positive signals.

В verifiable domains reward может быть назначен внешней процедурой, не зависящей от субъективного человеческого суждения. Для кода сгенерированную программу можно выполнить на test suite; для математики final answer можно распарсить и сравнить с canonical solution. В таких случаях событие ``answer is correct'' имеет смысл, внешний по отношению к самой language model. Современные reasoning models сильно опираются на эту структуру при обучении \citep{DeepSeek2025,Lightman2023}.

Сначала рассмотрим стандартный single-reference setting: для каждого prompt \(x\) существует ровно одна verifiable target \(y^{*}\). Определим reward
\[
    R(y,y^{*})=\mathbb{I}\{y=y^{*}\},
\]
где \(\mathbb{I}\{\cdot\}\) --- indicator function. Тогда objective равна
\[
    J_{\mathrm{RLVR}}(\theta\mid x,y^{*})
    =
    \mathbb{E}_{z\sim\pi_{\theta}(z\mid x)}
    \mathbb{E}_{y\sim\pi_{\theta}(y\mid x,z)}
    \big[R(y,y^{*})\big].
\]
Эквивалентно, в expanded summation form,
\[
    J_{\mathrm{RLVR}}(\theta\mid x,y^{*})
    =
    \sum_{z\in\mathcal{Z}}\sum_{y\in\mathcal{Y}}
    \pi_{\theta}(z\mid x)\pi_{\theta}(y\mid x,z)\,R(y,y^{*}).
\]

Для оптимизации этого expectation используется score-function identity, также известная как REINFORCE identity:
\[
    \nabla_{\theta}\,\mathbb{E}_{u\sim p_{\theta}(u)}[f(u)]
    =
    \mathbb{E}_{u\sim p_{\theta}(u)}
    \big[f(u)\nabla_{\theta}\log p_{\theta}(u)\big].
\]
Это тождество следует напрямую из соотношения \(\nabla_{\theta}p_{\theta}(u)=p_{\theta}(u)\nabla_{\theta}\log p_{\theta}(u)\). Применяя его к joint distribution of \((z,y)\), получаем
\begin{align*}
    \nabla_{\theta}J_{\mathrm{RLVR}}
     & =
    \nabla_{\theta}
    \sum_{z,y}\pi_{\theta}(z,y\mid x)\,R(y,y^{*})                                             \\
     & =
    \sum_{z,y}R(y,y^{*})\,\nabla_{\theta}\pi_{\theta}(z,y\mid x)                              \\
     & =
    \sum_{z,y}\pi_{\theta}(z,y\mid x)\,R(y,y^{*})\,\nabla_{\theta}\log\pi_{\theta}(z,y\mid x) \\
     & =
    \mathbb{E}_{z,y}\Big[
        R(y,y^{*})\,\nabla_{\theta}\log\pi_{\theta}(z,y\mid x)
        \Big].
\end{align*}
Используя факторизацию joint policy,
\[
    \log\pi_{\theta}(z,y\mid x)
    =
    \log\pi_{\theta}(z\mid x)+\log\pi_{\theta}(y\mid x,z),
\]
следовательно,
\[
    \nabla_{\theta}\log\pi_{\theta}(z,y\mid x)
    =
    \nabla_{\theta}\log\pi_{\theta}(z\mid x)
    +
    \nabla_{\theta}\log\pi_{\theta}(y\mid x,z).
\]
Поэтому
\[
    \nabla_{\theta}J_{\mathrm{RLVR}}
    =
    \mathbb{E}_{z,y}\Big[
        R(y,y^{*})
        \big(
        \nabla_{\theta}\log\pi_{\theta}(z\mid x)
        +
        \nabla_{\theta}\log\pi_{\theta}(y\mid x,z)
        \big)
        \Big].
\]

Это выражение имеет простую интерпретацию. Sampled pair \((z,y)\) получает reward только тогда, когда verifier объявляет \(y\) правильным. Когда это происходит, update увеличивает и probability reasoning trace, приведшего к ответу, и probability самого ответа при этом trace.

\subsection{VeriFree}

Юмор не предоставляет такого external verifier, и в более общем виде многие open-ended generation tasks не имеют deterministic correctness test. VeriFree закрывает этот разрыв, замечая, что в single-reference case expected exact-match reward можно переписать как conditional probability \citep{Zhou2025}.

Зафиксируем \(x\) и \(z\). Тогда
\begin{align*}
    \mathbb{E}_{y\sim\pi_{\theta}(\cdot\mid x,z)}
    \big[\mathbb{I}\{y=y^{*}\}\big]
     & =
    \sum_{y\in\mathcal{Y}}
    \pi_{\theta}(y\mid x,z)\,\mathbb{I}\{y=y^{*}\} \\
     & =
    \pi_{\theta}(y^{*}\mid x,z).
\end{align*}
Подставляя это тождество обратно в RLVR objective, получаем
\[
    J_{\mathrm{RLVR}}(\theta\mid x,y^{*})
    =
    \mathbb{E}_{z\sim\pi_{\theta}(z\mid x)}
    \big[\pi_{\theta}(y^{*}\mid x,z)\big].
\]
Это мотивирует verifier-free reward
\[
    r_{\theta}(x,z,y^{*}) := \pi_{\theta}(y^{*}\mid x,z),
\]
и соответствующую objective
\[
    J_{\mathrm{VF}}(\theta\mid x,y^{*})
    :=
    \mathbb{E}_{z\sim\pi_{\theta}(z\mid x)}
    \big[r_{\theta}(x,z,y^{*})\big].
\]
При single-reference assumption
\[
    J_{\mathrm{VF}}(\theta\mid x,y^{*}) \equiv J_{\mathrm{RLVR}}(\theta\mid x,y^{*})
\]
как scalar objectives. Значимость такой переформулировки одновременно практическая и концептуальная: для вычисления reward не требуется явное decoding \(y\), а \(\pi_{\theta}(y^{*}\mid x,z)\) можно вычислить teacher forcing на reference sequence, что удобно параллелизуется.

Чтобы получить policy gradient, запишем
\[
    J_{\mathrm{VF}}(\theta\mid x,y^{*})
    =
    \sum_{z\in\mathcal{Z}} \pi_{\theta}(z\mid x)\,r_{\theta}(x,z,y^{*}).
\]
Дифференцируя по product rule,
\begin{align*}
    \nabla_{\theta}J_{\mathrm{VF}}
     & =
    \sum_{z}
    \Big[
        \nabla_{\theta}\pi_{\theta}(z\mid x)\,r_{\theta}(x,z,y^{*})
        +
        \pi_{\theta}(z\mid x)\,\nabla_{\theta}r_{\theta}(x,z,y^{*})
    \Big] \\
     & =
    \sum_{z}\pi_{\theta}(z\mid x)
    \Big[
        r_{\theta}(x,z,y^{*})\,\nabla_{\theta}\log\pi_{\theta}(z\mid x)
        +
        \nabla_{\theta}r_{\theta}(x,z,y^{*})
        \Big].
\end{align*}
Так как \(r_{\theta}(x,z,y^{*})=\pi_{\theta}(y^{*}\mid x,z)\),
\[
    \nabla_{\theta}r_{\theta}(x,z,y^{*})
    =
    \pi_{\theta}(y^{*}\mid x,z)\,\nabla_{\theta}\log\pi_{\theta}(y^{*}\mid x,z)
    =
    r_{\theta}(x,z,y^{*})\,\nabla_{\theta}\log\pi_{\theta}(y^{*}\mid x,z).
\]
Следовательно,
\[
    \nabla_{\theta}J_{\mathrm{VF}}
    =
    \mathbb{E}_{z\sim\pi_{\theta}(z\mid x)}\Big[
        r_{\theta}(x,z,y^{*})
        \big(
        \nabla_{\theta}\log\pi_{\theta}(z\mid x)
        +
        \nabla_{\theta}\log\pi_{\theta}(y^{*}\mid x,z)
        \big)
        \Big].
\]

Эта декомпозиция информативна. Первый член,
\[
    \mathbb{E}_{z}\big[r_{\theta}(x,z,y^{*})\,\nabla_{\theta}\log\pi_{\theta}(z\mid x)\big],
\]
является score-function term по sampled reasoning traces. Он повышает вероятность traces, при которых сама модель назначает более высокую conditional mass reference answer. Второй член,
\[
    \mathbb{E}_{z}\big[\nabla_{\theta}r_{\theta}(x,z,y^{*})\big],
\]
является direct teacher-forced gradient на reference sequence. По сравнению с RLVR, VeriFree убирает шум от sampling final answer \(y\): estimator все еще samples \(z\), но не answer, используемый для reward evaluation. Именно поэтому VeriFree привлекателен в доменах, где verifier недоступен, но single reference answer существует.

\subsection{Multi-reference VeriFree}

Single-reference assumption не подходит для юмора. Prompt вроде ``write a joke using the keywords cat, dog'' не имеет одного правильного completion; он имеет много правдоподобных completions, различающихся setup, punchline mechanism, tone и уровнем абсурдности. Фактически реализованный репозиторий явно предназначен для построения нескольких референсов для одного weak prompt. Поэтому reward должен быть определен на множестве acceptable answers, а не на одном ответе.

Пусть \(Y^{*}(x)\subset\mathcal{Y}\) обозначает observed set of acceptable references для prompt \(x\). Определим multi-reference verified reward
\[
    R(y,Y^{*})=\mathbb{I}\{y\in Y^{*}\}.
\]
Соответствующая set-valued exact-match objective равна
\[
    J_{\mathrm{MR}}(\theta\mid x,Y^{*})
    =
    \mathbb{E}_{z\sim\pi_{\theta}(z\mid x)}
    \mathbb{E}_{y\sim\pi_{\theta}(\cdot\mid x,z)}
    \big[\mathbb{I}\{y\in Y^{*}\}\big].
\]
Для фиксированных \(x\) и \(z\),
\begin{align*}
    \mathbb{E}_{y\sim\pi_{\theta}(\cdot\mid x,z)}
    \big[\mathbb{I}\{y\in Y^{*}\}\big]
     & =
    \sum_{y\in\mathcal{Y}}
    \pi_{\theta}(y\mid x,z)\,\mathbb{I}\{y\in Y^{*}\} \\
     & =
    \sum_{y\in Y^{*}}
    \pi_{\theta}(y\mid x,z).
\end{align*}
Это мотивирует multi-reference verifier-free reward
\[
    r_{\theta}(x,z,Y^{*})
    :=
    \sum_{y\in Y^{*}}
    \pi_{\theta}(y\mid x,z).
\]
Величина \(r_{\theta}(x,z,Y^{*})\) лежит в \([0,1]\): это total probability mass, которую policy назначает, при reasoning trace \(z\), observed reference set. Эквивалентно, это вероятность события, что generated answer попадает внутрь этого множества. Итоговая objective равна
\[
    J_{\mathrm{MRVF}}(\theta\mid x,Y^{*})
    =
    \mathbb{E}_{z\sim\pi_{\theta}(z\mid x)}
    \big[r_{\theta}(x,z,Y^{*})\big].
\]

На scalar level MRVF играет для set-valued reward ту же роль, что VeriFree играет для single exact-match reward. Однако gradient structure не является просто single-reference формулой, где \(y^{*}\) заменено символом множества. Answer-side derivative меняется существенно.

Действительно,
\[
    J_{\mathrm{MRVF}}(\theta\mid x,Y^{*})
    =
    \sum_{z}\pi_{\theta}(z\mid x)\,r_{\theta}(x,z,Y^{*}),
\]
поэтому
\begin{align*}
    \nabla_{\theta}J_{\mathrm{MRVF}}
     & =
    \sum_{z}
    \Big[
        \nabla_{\theta}\pi_{\theta}(z\mid x)\,r_{\theta}(x,z,Y^{*})
        +
        \pi_{\theta}(z\mid x)\,\nabla_{\theta}r_{\theta}(x,z,Y^{*})
    \Big] \\
     & =
    \mathbb{E}_{z\sim\pi_{\theta}(z\mid x)}
    \Big[
        r_{\theta}(x,z,Y^{*})\,\nabla_{\theta}\log\pi_{\theta}(z\mid x)
        +
        \nabla_{\theta}r_{\theta}(x,z,Y^{*})
        \Big].
\end{align*}
Теперь раскроем reference-side derivative:
\begin{align*}
    \nabla_{\theta}r_{\theta}(x,z,Y^{*})
     & =
    \nabla_{\theta}
    \sum_{y\in Y^{*}}\pi_{\theta}(y\mid x,z) \\
     & =
    \sum_{y\in Y^{*}}
    \nabla_{\theta}\pi_{\theta}(y\mid x,z)   \\
     & =
    \sum_{y\in Y^{*}}
    \pi_{\theta}(y\mid x,z)\,\nabla_{\theta}\log\pi_{\theta}(y\mid x,z).
\end{align*}
Поэтому
\[
    \nabla_{\theta}J_{\mathrm{MRVF}}
    =
    \mathbb{E}_{z\sim\pi_{\theta}(z\mid x)}
    \Bigg[
        r_{\theta}(x,z,Y^{*})\,\nabla_{\theta}\log\pi_{\theta}(z\mid x)
        +
        \sum_{y\in Y^{*}}
        \pi_{\theta}(y\mid x,z)\,\nabla_{\theta}\log\pi_{\theta}(y\mid x,z)
        \Bigg].
\]

Иногда полезно переписать второй член в нормализованной форме. Если \(r_{\theta}(x,z,Y^{*})>0\), определим
\[
    w_{\theta}(y\mid x,z,Y^{*})
    :=
    \frac{\pi_{\theta}(y\mid x,z)}{r_{\theta}(x,z,Y^{*})}
    \,\mathbb{I}\{y\in Y^{*}\}.
\]
Тогда \(w_{\theta}(\cdot\mid x,z,Y^{*})\) является probability distribution over observed reference set, и
\[
    \nabla_{\theta}r_{\theta}(x,z,Y^{*})
    =
    r_{\theta}(x,z,Y^{*})
    \sum_{y\in Y^{*}}
    w_{\theta}(y\mid x,z,Y^{*})
    \nabla_{\theta}\log\pi_{\theta}(y\mid x,z).
\]
Эквивалентно,
\[
    \nabla_{\theta}\log r_{\theta}(x,z,Y^{*})
    =
    \sum_{y\in Y^{*}}
    w_{\theta}(y\mid x,z,Y^{*})
    \nabla_{\theta}\log\pi_{\theta}(y\mid x,z),
\]
и, следовательно,
\[
    \nabla_{\theta}J_{\mathrm{MRVF}}
    =
    \mathbb{E}_{z\sim\pi_{\theta}(z\mid x)}
    \Big[
        r_{\theta}(x,z,Y^{*})
        \big(
        \nabla_{\theta}\log\pi_{\theta}(z\mid x)
        +
        \nabla_{\theta}\log r_{\theta}(x,z,Y^{*})
        \big)
        \Big].
\]

Эта форма проясняет отличие от single-reference VeriFree. В VeriFree direct term --- это gradient \(\log\pi_{\theta}(y^{*}\mid x,z)\) для одного reference. В MRVF direct term --- это gradient \(\log \sum_{y\in Y^{*}}\pi_{\theta}(y\mid x,z)\). Поэтому update не отдает априорного предпочтения одному reference; он увеличивает total mass на observed set, причем каждый reference взвешивается текущей policy mass, назначенной ему. Если один reference уже доминирует, gradient концентрируется на нем. Если mass распределена по нескольким references, gradient усиливает их совместно. Поэтому MRVF не является просто нотационным вариантом VeriFree: optimization geometry меняется с одной target string на log-sum по многим targets.

\subsection{Неполные наблюдаемые референсы}

Observed set \(Y^{*}(x)\) все еще не является полным множеством хороших шуток. Возникают две разные формы incompleteness.

\paragraph{Structural incompleteness.}
Даже если репозиторий извлекает несколько references для prompt, может существовать много дополнительных шуток в latent set \(Y_{\mathrm{true}}(x)\), которые также валидны, но не наблюдаются. Это неизбежно в creative domains. Конечный датасет не может перечислить все acceptable punchlines для weak prompt.

\paragraph{Computational incompleteness.}
Даже внутри observed set \(Y^{*}(x)\) вычисление
\[
    r_{\theta}(x,z,Y^{*})=\sum_{y\in Y^{*}}\pi_{\theta}(y\mid x,z)
\]
может стать дорогим, когда \(|Y^{*}|\) велико. Поэтому сумму можно оценивать на sampled subset \(Y\subseteq Y^{*}\).

Предположим, что \(Y\) sampled uniformly without replacement из \(Y^{*}\), и каждый reference \(y\in Y^{*}\) имеет inclusion probability
\[
    \rho = P(y\in Y)=\frac{|Y|}{|Y^{*}|}.
\]
Определим subset reward
\[
    r_{\theta}(x,z,Y)
    :=
    \sum_{y\in Y}\pi_{\theta}(y\mid x,z)
    =
    \sum_{y\in Y^{*}}
    \mathbb{I}\{y\in Y\}\,\pi_{\theta}(y\mid x,z).
\]
Беря expectation по случайному subset,
\begin{align*}
    \mathbb{E}_{Y}\big[r_{\theta}(x,z,Y)\big]
     & =
    \sum_{y\in Y^{*}}
    \mathbb{E}_{Y}\big[\mathbb{I}\{y\in Y\}\big]\,
    \pi_{\theta}(y\mid x,z)                  \\
     & =
    \sum_{y\in Y^{*}}
    P(y\in Y)\,\pi_{\theta}(y\mid x,z)       \\
     & =
    \rho
    \sum_{y\in Y^{*}}\pi_{\theta}(y\mid x,z) \\
     & =
    \rho\,r_{\theta}(x,z,Y^{*}).
\end{align*}
Следовательно, subset reward пропорционален в expectation full observed-set reward.

Здесь становится релевантным GRPO baselining. Предположим, что для фиксированного prompt \(x\) все sampled rewards в группе умножаются на один и тот же positive factor \(\rho_x\), например потому что используется только доля observed references. Если стандартизовать rewards внутри группы,
\[
    \widetilde A_i
    =
    \frac{\hat r_i - \mu_x}{\sigma_x+\varepsilon},
    \qquad
    \mu_x=\frac{1}{G}\sum_{j=1}^{G}\hat r_j,
\]
то умножение каждого \(\hat r_i\) на \(\rho_x>0\) умножает и numerator, и denominator на один и тот же factor, оставляя \(\widetilde A_i\) неизменным. В этом смысле GRPO group standardization scale-invariant. Она смягчает reward-magnitude shifts в score-function term, хотя не восстанавливает информацию, потерянную из-за missing references, и сама по себе не решает bias direct reference gradient.

Другая фундаментальная проблема --- false-negative problem. Оптимизация
\[
    r_{\theta}(x,z,Y^{*})=\sum_{y\in Y^{*}}\pi_{\theta}(y\mid x,z)
\]
увеличивает probability mass на observed references. Поскольку conditional distribution по всем строкам суммируется в единицу,
\[
    \sum_{y\in\mathcal{Y}}\pi_{\theta}(y\mid x,z)=1,
\]
увеличенная mass должна прийти из \(\mathcal{Y}\setminus Y^{*}\). К сожалению, это complement содержит не только плохие шутки, но и ненаблюдаемые хорошие шутки из \(Y_{\mathrm{true}}(x)\setminus Y^{*}(x)\). Иными словами, training signal может непреднамеренно подавлять валидные, но unseen humorous continuations. Этот риск не является дефектом gradient derivation; он является структурным следствием learning from incomplete supervision в intrinsically many-solution domain.

Чтобы уменьшить этот риск, предлагаемая оптимизация включает Kullback-Leibler penalty относительно reference policy \(\pi_{\mathrm{ref}}\):
\[
    \max_{\theta}
    \;
    \mathbb{E}_{x,z}\big[r_{\theta}(x,z,Y^{*}(x))\big]
    -
    \beta\,
    \mathbb{E}_{x}\Big[
        \mathrm{KL}\big(
        \pi_{\theta}(\cdot\mid x)\,\|\,\pi_{\mathrm{ref}}(\cdot\mid x)
        \big)
        \Big].
\]
Эквивалентно, можно рассматривать задачу в constrained form,
\[
    \max_{\theta}\;
    \mathbb{E}_{x,z}\big[r_{\theta}(x,z,Y^{*}(x))\big]
    \quad
    \text{s.t.}
    \quad
    \mathbb{E}_{x}
    \Big[
        \mathrm{KL}\big(
        \pi_{\theta}(\cdot\mid x)\,\|\,\pi_{\mathrm{ref}}(\cdot\mid x)
        \big)
        \Big]
    \le \varepsilon.
\]
В данном контексте KL term не является просто generic regularization. Его методологическая роль --- препятствовать чрезмерно агрессивному перераспределению probability mass от reference model, когда observed reference sets заведомо неполны.

\subsection{Предлагаемый training loop}

Предполагаемый training stage samples несколько reasoning traces для одного prompt и использует их rewards для обновления reasoning policy. Для prompt \(x\) sample group
\[
    z_1,\dots,z_G \sim \pi_{\theta}(z\mid x).
\]
Для каждого \(z_i\) вычисляется multi-reference reward
\[
    \hat r_i
    :=
    r_{\theta}(x,z_i,Y^{*}(x))
    =
    \sum_{y\in Y^{*}(x)}
    \pi_{\theta}(y\mid x,z_i).
\]
На практике каждая \(\pi_{\theta}(y\mid x,z_i)\) получается teacher forcing на reference sequence \(y\). Поскольку эти sequence probabilities очень малы, реализация должна аккумулировать token log-probabilities и, где необходимо, использовать stabilization вроде \(\log\sum\exp\) перед обратным переходом к probabilities или normalized weights.

Gradient естественно раскладывается на три части:
\[
    \nabla_{\theta}J_{\mathrm{MRVF}}
    =
    g_{z}+g_{r}+g_{\mathrm{KL}}.
\]

\paragraph{Reasoning term.}
Score-function часть по \(z\) оценивается по sampled traces и поэтому выигрывает от variance reduction. Чистый выбор --- leave-one-out baseline
\[
    b_i
    =
    \frac{1}{G-1}\sum_{j\neq i}\hat r_j,
    \qquad
    A_i=\hat r_i-b_i.
\]
Соответствующий estimator равен
\[
    g_{z}
    =
    \frac{1}{G}\sum_{i=1}^{G}
    A_i\,\nabla_{\theta}\log\pi_{\theta}(z_i\mid x).
\]
Это та же логика, которая мотивирует RLOO estimators: каждый trace сравнивается с другими traces, sampled для того же prompt. При желании можно дополнительно нормализовать \(A_i\) внутри группы, чтобы уменьшить prompt-specific reward-scale variation, в духе GRPO. Такая нормализация повышает robustness, но уже не является строго unbiased, поскольку scale estimate сама зависит от sampled group.

\paragraph{Reference term.}
Для direct answer-side contribution корректный observed-set gradient равен
\[
    g_{r}
    =
    \frac{1}{G}\sum_{i=1}^{G}
    \nabla_{\theta}r_{\theta}(x,z_i,Y^{*}(x))
    =
    \frac{1}{G}\sum_{i=1}^{G}
    \sum_{y\in Y^{*}(x)}
    \pi_{\theta}(y\mid x,z_i)\,
    \nabla_{\theta}\log\pi_{\theta}(y\mid x,z_i).
\]
Эквивалентно, используя normalized within-set weights
\[
    w_{i,y}
    :=
    \frac{\pi_{\theta}(y\mid x,z_i)}
    {\hat r_i}
    \,\mathbb{I}\{y\in Y^{*}(x)\},
\]
можно записать
\[
    g_{r}
    =
    \frac{1}{G}\sum_{i=1}^{G}
    \hat r_i
    \sum_{y\in Y^{*}(x)}
    w_{i,y}\,
    \nabla_{\theta}\log\pi_{\theta}(y\mid x,z_i),
\]
если \(\hat r_i>0\). Эта форма делает geometry явной: каждый trace получает total weight \(\hat r_i\), а внутри этого trace gradient распределяется по references согласно текущей policy mass на каждом reference.

\paragraph{KL term.}
Наконец,
\[
    g_{\mathrm{KL}}
    =
    -\beta\,
    \nabla_{\theta}
    \mathrm{KL}\big(
    \pi_{\theta}(\cdot\mid x)\,\|\,\pi_{\mathrm{ref}}(\cdot\mid x)
    \big).
\]
Полный update для minibatch усредняет эти prompt-level contributions по sampled prompts.

Практическая важность реализованного пайплайна теперь ясна. Без prompt-conditioned multi-reference sets \(Y^{*}(x)\) нельзя вычислить ни \(\hat r_i\), ни \(g_r\). Поэтому data pipeline не является вспомогательным к теории; он предоставляет observed structure, от которой зависит MRVF optimization.

\subsection{Замечание о baselining}

Variance reduction прямолинейно применима только к score-function term. Причина в том, что \(z\) действительно sampled from \(\pi_{\theta}(z\mid x)\). Если \(b(x,z_{-i})\) --- baseline, не зависящий от sampled variable \(z_i\), то
\begin{align*}
    \mathbb{E}_{z_i\sim\pi_{\theta}(\cdot\mid x)}
    \big[
        b(x,z_{-i})\,\nabla_{\theta}\log\pi_{\theta}(z_i\mid x)
        \big]
     & =
    b(x,z_{-i})
    \sum_{z_i}\pi_{\theta}(z_i\mid x)\,
    \nabla_{\theta}\log\pi_{\theta}(z_i\mid x)       \\
     & =
    b(x,z_{-i})
    \sum_{z_i}\nabla_{\theta}\pi_{\theta}(z_i\mid x) \\
     & =
    b(x,z_{-i})\,\nabla_{\theta}1                    \\
     & = 0.
\end{align*}
Поэтому leave-one-out baseline unbiased для reasoning term.

Reference term устроен иначе. Это не expectation по sampled answers \(y\sim\pi_{\theta}(\cdot\mid x,z)\), а direct derivative summed reference probability. Следовательно, нет тождества вида
\[
    \mathbb{E}_{y}\big[\nabla_{\theta}\log\pi_{\theta}(y\mid x,z)\big]=0
\]
доступного для centering. Например, в single-reference case,
\[
    \mathbb{E}_{z}
    \big[
        (r_{\theta}(x,z,y^{*})-b(x))
        \nabla_{\theta}\log\pi_{\theta}(y^{*}\mid x,z)
        \big]
\]
равно
\[
    \mathbb{E}_{z}
    \big[
        r_{\theta}(x,z,y^{*})
        \nabla_{\theta}\log\pi_{\theta}(y^{*}\mid x,z)
        \big]
    -
    b(x)\,
    \mathbb{E}_{z}
    \big[
        \nabla_{\theta}\log\pi_{\theta}(y^{*}\mid x,z)
        \big],
\]
и второй член в общем случае не равен нулю. Та же проблема сохраняется в multi-reference case с \(\nabla_{\theta}\log r_{\theta}(x,z,Y^{*})\). Поэтому baselines, безвредные для score-function term, становятся biased, если наивно применить их к direct reference gradient.

Это различие важно для дизайна training algorithm. Reasoning term должен использовать variance-reduction machinery, такую как leave-one-out baselines или аккуратно выбранную group normalization. Reference term, напротив, следует вычислять максимально прямо из observed-set probability, принимая, что это deterministic, но потенциально biased proxy для true reward, порождаемого \(Y_{\mathrm{true}}(x)\).

В совокупности эти наблюдения объясняют общую структуру предложения. MRVF наследует ключевое вычислительное преимущество verifier-free learning: он избегает external verifiers и explicit reward sampling over final answers, но должен быть адаптирован к multi-reference и incomplete характеру юмора. Итоговая objective математически tractable, но ее ограничения являются частью вклада работы: incompleteness, false negatives и estimator design центральны для проблемы.
